\documentclass[10pt,a4paper]{amsart}
\usepackage[margin=19mm]{geometry}
\usepackage{amsmath,amssymb,mathtools}
\usepackage[scr=rsfs,cal=euler]{mathalfa}
\usepackage{xfrac}
\usepackage[unicode,hidelinks,bookmarks=false]{hyperref}
\newcommand{\ie}{i.e.\ }
\newcommand{\cf}{cf.\ }
\newcommand{\resp}{respectively\ }
\newcommand{\Subsection}{\subsection}
\providecommand{\nobreakdash}{\nobreak\hbox{-}}
\setlength{\emergencystretch}{2em}
\multlinegap0pt
\usepackage{bm}
\usepackage{bbm}
\usepackage{dsfont}
\usepackage{xspace}
\usepackage{cancel}
\usepackage{mathtools}
\usepackage{xcolor}
\usepackage{amsthm}
\makeatletter
\@ifundefined{theorem}{\newtheorem{theorem}{Theorem}}{}
\@ifundefined{lemma}{\newtheorem{lemma}[theorem]{Lemma}}{}
\makeatother
\newcommand{\GG}{\mathbb{G}}
\newcommand{\PP}{\mathbb P}
\newcommand{\f}{\varphi}
\newcommand{\kk}{\mathrm{k}}
\newcommand{\la}{\langle}
\newcommand{\ra}{\rangle}
\renewcommand{\L}{\Lambda}
\renewcommand{\a}{\alpha}
\renewcommand{\b}{\beta}
\title[8-dimensional 2-step nilpotent Lie algebras over algebraical]{8-dimensional 2-step nilpotent Lie algebras over algebraically closed fields of char $\ne 2, 3$}
\author{Giovanni Bazzoni, Juan Rojo}
\date{Source: arXiv:2507.14907v3 (CC BY 4.0)}
\begin{document}
\maketitle

\noindent\emph{Source and licence.} 8-dimensional 2-step nilpotent Lie algebras over algebraically closed fields of char $\ne 2, 3$. Version arXiv:2507.14907v3 (CC BY 4.0), distributed under the Creative Commons Attribution 4.0 International licence, as recorded in the arXiv licence metadata for this exact version and shown on the abstract page https://arxiv.org/abs/2507.14907v3. Full rights notice: licenses/arxiv-2507.14907-CC-BY-4.0.txt.\par\smallskip

\noindent\textit{Editorial scope.} Counted unit: the Lemma whose source label is \texttt{lem:rational-map} in the article (source0.tex lines 1516--1535, source label \texttt{lem:rational-map}), delivered together with the article's own complete original proof of it (source0.tex lines 1536--1608, one \texttt{proof} environment of 73 lines). No mathematics is written, repaired or re-derived: statement and proof are the article's own text line for line. Required context retained uncounted: lem:case16-not-same-line (lines 1473--1477). Every definition, display and label of those lines is retained unchanged. Omitted from this package and not redistributed: 1--1472, 1478--1515, 1609--2540. Nothing inside a retained excerpt is omitted or altered: every retained body is delivered exactly as the article's lines are written, so no omission receipt of any kind accompanies this package. Citations: the retained excerpts carry no citation command at all, so no bibliography entry of the article is transcribed and no reference is redistributed. Reference adaptation: none was needed and none was made; every reference inside the delivered unit resolves inside it, so no unresolved reference or citation is produced. Printed slips kept literal: no printed word, symbol, hypothesis or number of the counted statement or of its proof was altered. Layout and class: the article's own class is \texttt{amsart} with packages including latexsym, amsmath, amsfonts, amscd, amssymb, color, mathtools, graphicx, wrapfig, xy, hyperref, tabu, soul, mathrsfs; the wrapper is the lane's pinned \texttt{amsart} editorial wrapper at 10pt on A4, which restates the theorem-like environments the retained text needs and adds the article's own macro definitions that the retained text needs (\texttt{\textbackslash GG}, \texttt{\textbackslash L}, \texttt{\textbackslash PP}, \texttt{\textbackslash a}, \texttt{\textbackslash b}, \texttt{\textbackslash f}, \texttt{\textbackslash kk}, \texttt{\textbackslash la}, \texttt{\textbackslash ra}), with extra wrapper packages (bm, bbm, dsfont, xspace, cancel, mathtools, xcolor). The delivered numbering is the unit's own. Assets: no retained span contains a figure environment, a graphics-inclusion command, a PDF-page inclusion, a table or a TikZ picture; no image file is copied into this package, the package declares no asset dependency and no image file is redistributed with it. Licence: version arXiv:2507.14907v3 of this article is distributed under the Creative Commons Attribution 4.0 International licence, as recorded in the arXiv licence metadata for this exact version and shown on the abstract page https://arxiv.org/abs/2507.14907v3. A full rights notice is delivered at licenses/arxiv-2507.14907-CC-BY-4.0.txt. Characters: every retained excerpt is pure ASCII, so the delivered engine, XeLaTeX with its default Latin Modern text font, reports no missing character.
\begin{lemma} \label{lem:case16-not-same-line} 
Suppose $\pi =\la \f_1, \f_2, \f_3 \ra \subset \PP(\L^2 W)=\PP^9$ is a plane such that $\pi \cap \GG=\{[\f_1]\}$. Assume also that dim~$(U_{\f} \cap U_1)=1$ for any $\f \in \pi \setminus \{[\f_1]\}$. Then the lines $U_{\f} \cap U_1$ are not all the same. In other words, $\bigcap_{\f \in \pi} U_{\f}=\{0\}$. 

\noindent In particular, we can choose generators $\f_1, \f_2, \f_3$ so that $U_1=\la x_1, x_2 \ra$, $U_2 \cap U_1=\la x_1 \ra$, and $U_3 \cap U_1=\la x_2 \ra$.
\end{lemma}

\begin{lemma} \label{lem:rational-map}
Suppose $\pi =\la \f_1, \f_2, \f_3 \ra$ is a plane as in Lemma \ref{lem:case16-not-same-line}. Then, the map $\PP^1 \to \PP(U_1) \cong \PP^1$ given by
\[
[\a:\b] \mapsto U_{\a\f_2+\b\f_3} \cap U_1
\]
is a rational map of degree 2. With respect to a suitable choice of generators $\f_2, \f_3$ and basis $x_1,x_2$ of $U_1$, the map is given by the matrix:
\[
[\a : \b] \mapsto 
\begin{pmatrix}
    1 & a & 0 \\
    0 & b & 1
\end{pmatrix} 
\begin{pmatrix}
    \a^2 \\
    \a \b \\
    \b^2
\end{pmatrix}
\]
with $ab \ne 1$.
\end{lemma}

\begin{proof}
By the previous lemma we know the above map is non-constant. We can choose a basis $\{x_i\}$ and generators $\f_1, \f_2, \f_3$ for $\pi$ so that $U_1=\la x_1, x_2 \ra$, $U_2=\la x_1, x_3,x_4,x_5 \ra$ and $\la x_2 \ra= U_1 \cap U_3$. Then 
\[
\f_2=x_1(ax_3+bx_4+cx_5)+x_3(ex_4+fx_5)+gx_4x_5\,,
\]
and at least one of $a, b,c$ is non-zero. We can assume (swapping coordinates maybe) that $a \ne 0$, and rescaling $x_3$ we get $a=1$. Make the change $x'_3=x_3+bx_4+cx_5$, then $\f_2=x_1x'_3+x'_3(ex_4+fx_5)+g'x_4x_5$, with $g' \ne 0$, so rescaling $x_5$ we get $g'=1$. Reset notation, so $\f_2=x_1x_3+x_4(x_5-ex_3)+fx_3x_5$. Make the change $x'_5=x_5-ex_3$ so $\f_2=x_1x_3+x_4x'_5+fx_3x'_5$, reset notation so $\f_2=x_1x_3+(x_4+fx_3)x_5$, and put $x'_4=x_4+fx_3$. 

\item We start with $\f_1=x_1x_2$, $\f_2=x_1x_3+x_4x_5$ and $U_1\cap U_3=\la x_2 \ra$. For $j=3,4,5$ the affine line $x_j + \la x_1 \ra$ intersects $U_3$ in a point, hence we find $a_j \in \kk$ so that $x_j+a_jx_1 \in U_3$. In other words, through the change $x'_j=x_j+a_jx_1$, we can assume that $x_j \in U_3$, for $j=3,4,5$. Make these changes, and $\f_2=x_1x'_3+(x'_4-a_4x_1)(x'_5-a_5x_1)=x_1(x'_3+a_5x'_4\textcolor{red}{-}a_4x'_5)+x'_4x'_5$, so by a further change $x''_3=x'_3+a_5x'_4-a_4x'_5 \in U_3$, followed by relabeling, we get $\f_2=x_1x_3+x_4x_5$ and $U_3=\la x_2,x_3,x_4, x_5 \ra$. We have arrived to the following:
\[
\begin{cases}
    \f_1=x_1x_2 \\
    \f_2=x_1x_3+x_4x_5 \\
    \f_3=x_2(ax_3+bx_4+cx_5)+x_3(ex_4+fx_5)+gx_4x_5
\end{cases}
\]
If $b=c=0$ then $\f_3=(ax_2-dx_4-ex_5)x_3+fx_4x_5$ and $\f_3-f\f_2$ would have rank $2$, a contradiction. Hence $b$ or $c \ne 0$. Swapping the coordinates $x_4$ and $x_5$ (and maybe changing the sign of $\f_2$) we can assume that $c \ne 0$, so rescaling $x'_5=cx_5$, $x'_4=\frac1c x_4$ we can assume that $c=1$, and make the change $x'_5=x_5+bx_4$, so we can assume that $\f_3=x_2(ax_3+x_5)+x_3(ex_4+fx_5)+gx_4x_5$.

\item We claim now that $e\ne 0$. If it were $e=0$ then $\f_3=ax_2x_3+(x_2+fx_3+gx_4)x_5$, with $g \ne 0$, for otherwise the rank would drop. Then, the 4-space associated to the bivector $g\f_2-\f_3$ contains $U_1$, a contradiction.
%\textcolor{red}{(since otherwise $x_4 \notin U_{\f_3}$). A simple computation shows that
%\[
%U_{\a \b}\coloneqq U_{\a \f_2+\b \f_3}=
%\begin{cases}
%    \a x_3 \\
%    a \b x_3+\b x_5 \\
%    a \b x_2+\a x_1-\b e x_5 \\
%    (\a+f\b)x_5 \\
%    (\a+f\b)x_4+\b x_2+\b e x_3
%\end{cases}
%\]
%\[
%U_{\a \b}\coloneqq U_{\a \f_2+\b \f_3}=\la\a x_3,a \b x_3+\b x_5,a \b x_2+\a x_1-\b f x_5,(\a+g\b)x_5,(\a+g\b)x_4+\b x_2+\b f x_3\ra
%\]
%For $a\b\ne 0$ this space is generated by $x_3, x_5, a \b x_2+\a x_1, (\a+g \b)x_4+\b x_2$. If we choose $(\a_0:\b_0)$ with $\a_0 + g \b_0=0$, i.e. $\b_0=1, \a_0=-f \ne 0$, then $x_1, x_2 \in U_{\a_0 \b_0}$, a contradiction.} 
We conclude that $e \ne 0$, so we can assume $e=1$ by rescaling $x'_4=ex_4$, $x'_1=ex_1$, and we get 
\[
\f_3=x_2(ax_3+x_5)+x_3(x_4+fx_5)+gx_4x_5=x_2(ax_3+x_5)-x_3x'_4+gx'_4x_5
\]
with a further change $x'_4=-(x_4+fx_5)$, $x'_5=-x_5$. Hence we arrive at a considerably simplified model
\begin{equation} \label{eq:model16-preliminary}
 \begin{cases}
    \f_1=x_1x_2 \\
    \f_2=x_1x_3+x_4x_5 \\
    \f_3=(ax_2+x_4)x_3+(x_2+bx_4)x_5
\end{cases}   
\end{equation}
where we have relabeled the constants $a,b$, with $ab \ne 1$ since $\f_3$ has rank $4$. Let us compute the map $[\a:\b] \mapsto U_{\a \b } \cap U_1$, with $U_{\a \b}\coloneqq U_{\a \f_2+ \b \f_3}$. We have $\a \f_2+\b \f_3=(\a x_1+a\b x_2+\b x_4)x_3+(\b x_2+(\a+b\b)x_4)x_5$, so it follows easily that
\[
U_{\a \b}=\la\a x_3,\b ax_3+\b x_5,\a x_1+\b a x_2+\b x_4,\b x_3 + (\a + b \b) x_5,\b x_2+(\a+b \b) x_4\ra
\]
and for $\a\b\ne 0$ this is generated by $x_3$,$x_5$, $y_1=\a x_1+\b ax_2+\b x_4$, and $y_2=\b x_2+(\a+b \b)x_4$. We eliminate $x_4$ by considering 
\[
(\a +b \b)y_1-\b y_2=(\a^2+b \a \b)x_1+(\tfrac{a}{ab-1} \a \b +\b^2) (ab-1)x_2
\]
with $ab-1 \ne 0$. With respect to the basis $x'_1=x_1$, and $x'_2=(ab-1)x_2$ of $U_1$, we have obtained the degree two rational map 
\[
(\a: \b) \mapsto 
\begin{pmatrix}
    \a^2+b\a \b \\
    \tfrac{a}{ab-1} \a \b+\b^2
\end{pmatrix}
=
\begin{pmatrix}
    1 & b & 0 \\
    0 & \tfrac{a}{ab-1} & 1
\end{pmatrix} 
\begin{pmatrix}
    \a^2 \\
    \a \b \\
    \b^2
\end{pmatrix}
\]
with $\tfrac{ab}{ab-1} \ne 1$, as desired.
\end{proof}

\end{document}
